\section{The limit density at fixed memory}\label{sec:bulk-applications}

Fix $\varepsilon\in(0,\frac14)$, so $a=1-2\varepsilon\in(\frac12,1)$.
Conditional on $X_1=1$, $S_k/k^a$ converges almost surely to
$\mathcal L_a$~\cite[Theorem~3.7]{bercu2018}. The fair-start limit
$\mathcal L^\circ_a$ has the law of $\xi\mathcal L_a$, where $\xi$ is a fair sign independent of
$\mathcal L_a$~\cite[(1.3) and (1.6)]{glr2026}. Gu\'erin, Laulin and Raschel
prove that both limits have positive, bounded and smooth
densities~\cite[Theorem~1.3]{glr2026}, denoted by $h_a$ and $g_a$ as before.
Gu\'erin, Laulin, Raschel and Simon prove
that $\mathcal L_a$ is unimodal~\cite[Theorem~1.1]{glrs2025}.

\subsection{Memory near one}

\begin{corollary}[the fixed-memory limit near one]\label{cor:fixed-memory-bridge}
For the elephant walk with $X_1=1$ and fixed memory $p=(1+a)/2$,
$a\in(\frac12,1)$, let $\mathcal L_a=\lim_{k\to\infty}S_k/k^a$
as in~\cite[Theorem~3.7]{bercu2018}. If $Z$ has characteristic function
$\mathbb E e^{itZ}=\exp\{-it\log|t|-\pi|t|/2\}$, continuously extended at
$t=0$, then
\[
 \frac{1-\mathcal L_{1-2\varepsilon}}{2\varepsilon}
       +\log(1/\varepsilon)\ \Longrightarrow\ Z
       \qquad(\varepsilon\downarrow0).
\]
\end{corollary}

\begin{proof}
Fix $0<\varepsilon<1/8$, put $a=1-2\varepsilon$, and use one infinite
walk with $X_1=1$. By~\cite[(2.4), Appendix~A]{bercu2018},
$\mathcal M_k=\Gamma(k)S_k/\Gamma(k+a)$ is a martingale with limit
$\mathcal L_a$. It is Bercu's martingale divided by $\Gamma(1+a)$.
Its increments are
\[
 \mathcal M_{k+1}-\mathcal M_k
 =\frac{\Gamma(k+1)}{\Gamma(k+1+a)}
       \left(X_{k+1}-\frac{aS_k}{k}\right).
\]
Log-convexity of gamma gives
$1\le k^a\Gamma(k)/\Gamma(k+a)\le(1+a/k)^{1-a}\le1+1/k$.
Since $S_k=k-2M_k$, Lemma~\ref{lem:B-tail} gives
\[
 \mathbb E\!\left[1-a^2(S_k/k)^2\right]
 \le4\varepsilon+\frac{4\mathbb EM_k}{k}
 \le4\varepsilon(1+\log k).
\]
Orthogonality of the increments therefore gives square-mean convergence and
\[
 \mathbb E|\mathcal L_a-\mathcal M_n|^2
 \le C\varepsilon\sum_{k\ge n}k^{-2a}(1+\log k)
 \le C\varepsilon n^{1-2a}(1+\log n),
\]
with $C$ independent of $a\in[3/4,1)$, since $2a-1\ge1/2$.

Also $|\mathcal M_n-n^{-a}S_n|\le n^{-a}$.
Appendix~\ref{app:fixed-refinements} combines these two errors with
the compound Poisson approximation and chooses $n$ as a function of
$\varepsilon$. The resulting Kolmogorov bound
\eqref{eq:near-one-kolmogorov} tends to zero, proving the claim.
\end{proof}

Theorem~\ref{thm:near-one-profile} proves the stronger uniform density
approximation and total variation convergence. Trinas and Valle find
$\Var\mathcal L_a=2(1-a)+o(1-a)$~\cite[Eq.~(4.5)]{tv2026}. In both results, $\mathcal L_a$ concentrates at $1$ as $a\to1$.
For Theorem~\ref{thm:two-modes-near-one}, however, the exact second
and fourth moments are enough.

\subsection{The number of modes}\label{sec:modes}

A probability density $h$ on $\mathbb R$ is \emph{unimodal} if for some $m$
it is nondecreasing on $(-\infty,m]$ and nonincreasing on $[m,\infty)$.
Gu\'erin, Laulin, Raschel and Simon ask for a unique local maximum~\cite{glrs2025},
which also rules out a flat top. The density $g_a$ is a mixture of the unimodal density of
$\mathcal L_a$ and its reflection, but mixing unimodal densities can
give more than one peak~\cite{dj1988}. First take a fair start.
Lemma~\ref{lem:analytic-density} rules out constant intervals, so
the modes are points. We consider other biases afterward.

\begin{theorem}[at most two modes]\label{thm:two-modes}
Let $a\in(\frac12,1)$. Exactly one of the following holds.
\begin{enumerate}[(a)]
\item The density $g_a$ is nondecreasing on $(-\infty,0]$ and nonincreasing on
$[0,\infty)$.
\item There is $\kappa>0$ such that $g_a$ is nondecreasing on $[0,\kappa]$ and
nonincreasing on $[\kappa,\infty)$, and $g_a(0)<g_a(\kappa)$.
\end{enumerate}
In case (b), $g_a$ is nondecreasing on $(-\infty,-\kappa]$ and nonincreasing on
$[-\kappa,0]$, and it is not unimodal. All these monotonicities are strict.
The unique local maximum is $0$ in case (a), and the only local maxima
are $\pm \kappa$ in case (b).
In case (a), there is $d>0$ such that $g_a$ is strictly concave on
$(-d,d)$ and strictly convex on each tail. In case (b), there are
$0<d_1<\kappa<d_2$ such that $g_a$ is strictly convex on $(-d_1,d_1)$ and
on each tail beyond $\pm d_2$, and strictly concave on
$(-d_2,-d_1)$ and $(d_1,d_2)$. Thus it has exactly two inflection points
in case (a) and four in case (b).
\end{theorem}

We cannot pass directly from the second differences of $P_N$ to
the curvature of $g_a$. First smooth with a Gaussian and apply
Descartes' rule to the resulting exponential sum. Now let
$N\to\infty$, keeping the width fixed, and only then send the width
to $0$. Appendix~\ref{app:modes} proves the lemmas used below. A real
function $v$ has \emph{values at most} $\Pi$ if
$\operatorname{sp}(v(y_1),\dots,v(y_j))\preceq\Pi$ for all $y_1<\dots<y_j$, and
\emph{increments at most} $\Pi$ if
$\operatorname{sp}(v(y_2)-v(y_1),\dots,v(y_j)-v(y_{j-1}))\preceq\Pi$ for all
such points. We put $\Pi_4=(+,-,+,-)$.

\begin{proof}
Fix $a$, so that $\varepsilon=(1-a)/2\in(0,\frac14)$.

\emph{Step 1, smoothing.} Fix $\sigma>0$ and let
$\varphi(y)=e^{-y^2/(2\sigma^2)}/(\sigma\sqrt{2\pi})$. For $N\ge3$ put
$\eta=2N^{-a}$ and $z_k=(2k-N)N^{-a}$, so that $S_N/N^a=z_{A_N}$, and let
\[
 V_N(y)=\eta^{-2}\sum_{k=1}^{N-1}D_N(k)\,\varphi(y-z_k).
\]
Since $\varphi(y-z_k)=\varphi(y)e^{-z_k^2/(2\sigma^2)}e^{yz_k/\sigma^2}$ and
the $z_k$ increase with $k$, Lemma~\ref{lem:descartes} and
Theorem~\ref{thm:five-blocks} show that $V_N$ has values at most $\Pi_5$. (If
every $D_N(k)$ vanishes, then $V_N=0$ and every pattern is empty.)

\emph{Step 2, the limit $N\to\infty$.} Let $u(y)=\E\varphi(y-\mathcal L^\circ_a)$.
The derivatives of $\varphi$ are bounded, so $u$ is smooth and
$u^{(j)}(y)=\E\varphi^{(j)}(y-\mathcal L^\circ_a)$. We show that
$V_N(y)\to u''(y)$ for every $y$. Put $P_N(k)=0$ for $k\notin\{0,\dots,N\}$.
Then $D_N(k)=0$ unless $-1\le k\le N+1$, and summation by parts gives
\[
 \sum_{k=-1}^{N+1}D_N(k)\,\varphi(y-z_k)=\sum_{k=0}^NP_N(k)\,\Delta\varphi(y-z_k),
 \qquad \Delta\varphi(t)=\varphi(t+\eta)-2\varphi(t)+\varphi(t-\eta).
\]
Taylor's formula gives
$\lvert\eta^{-2}\Delta\varphi(t)-\varphi''(t)\rvert\le\eta^2\lVert\varphi''''\rVert_\infty/12$.
By symmetry the four terms with $k\in\{-1,0,N,N+1\}$ on the left are at most
$\varphi(0)\,(6P_N(0)+2P_N(1))$ in absolute value. Here
$P_N(0)\le(1-\varepsilon)^{N-1}$, and the proof of Theorem~\ref{thm:dip} gives
$f_N(1)\le\varepsilon N(1-\varepsilon)^{N-2}$ and
$f_N(N-1)\le\varepsilon(1-\varepsilon)^{N-2}$, so that
$P_N(1)\le\varepsilon N(1-\varepsilon)^{N-2}$. As $\eta^{-2}=N^{2a}/4$,
these four terms times $\eta^{-2}$ tend to $0$. Hence
\[
 V_N(y)=\E\varphi''\bigl(y-S_N/N^a\bigr)+o(1)\longrightarrow
 \E\varphi''(y-\mathcal L^\circ_a)=u''(y),
\]
because $S_N/N^a\to\mathcal L^\circ_a$ in law and $\varphi''$ is bounded and
continuous. By Lemma~\ref{lem:sp-limit}, $u''$ has values at most $\Pi_5$.

\emph{Step 3, integrating the sign bound.} By dominated convergence, $u'(y)\to0$ as
$y\to\pm\infty$. So $u$ has increments at most $\Pi_4$, by
Lemma~\ref{lem:sp-derivative}.

\emph{Step 4, the limit $\sigma\to0$.} Write $u=u_\sigma$. As $g_a$ is bounded
and continuous, $u_\sigma(y)=\int g_a(y-t)\varphi(t)\,dt$ tends to $g_a(y)$ for
every $y$. So the increments of $u_\sigma$ tend to those of $g_a$, and
Lemma~\ref{lem:sp-limit} shows that $g_a$ has increments at most $\Pi_4$.
Lemma~\ref{lem:analytic-density} also makes $g_a''$ bounded and continuous,
so $u_\sigma''=g_a''*\varphi\to g_a''$ pointwise. Hence $g_a''$ has
values at most $\Pi_5$.

\emph{Step 5, symmetry.} Let $0\le y_1<\dots<y_j$, let
$d_i=g_a(y_{i+1})-g_a(y_i)$, and let $p$ be the pattern of the $d_i$. As $g_a$
is even, the points $-y_j<\dots<-y_1<y_1<\dots<y_j$, with $-y_1$ left out if
$y_1=0$, have the increments $-d_{j-1},\dots,-d_1$, then $0$ if $y_1>0$, then
$d_1,\dots,d_{j-1}$. Their pattern is the reverse of $-p$ followed by $p$. It
has $2\lvert p\rvert$ entries and is $\preceq\Pi_4$. So $p$ is empty, $(+)$,
$(-)$ or $(+,-)$. In particular there are no $0\le y<y'<y''$ with
$g_a(y)>g_a(y')<g_a(y'')$.

If no increment on $[0,\infty)$ is positive, (a) holds by evenness. Otherwise
$g_a(y_1)<g_a(y_2)$ for some $0\le y_1<y_2$. Some increment on
$[y_2,\infty)$ is negative, since otherwise $g_a\ge g_a(y_2)>0$ there, against
$\int g_a=1$. If $g_a(y_4)<g_a(y_3)$ with $y_2\le y_3<y_4$, then $g_a$ cannot rise
after $y_4$, so it attains its maximum over $[0,\infty)$ at some
$\kappa\in[0,y_4]$. If $0\le y<y'\le \kappa$ and $g_a(y)>g_a(y')$, then $g_a(y')<g_a(\kappa)$,
and $(y,y',\kappa)$ is a fall followed by a rise. If $\kappa\le y<y'$ and
$g_a(y)<g_a(y')$, then $(\kappa,y,y')$ is one. So $g_a$ is nondecreasing on $[0,\kappa]$
and nonincreasing on $[\kappa,\infty)$. If $g_a(0)=g_a(\kappa)$, then $g_a$ would be
nonincreasing on $[0,\infty)$, against $g_a(y_1)<g_a(y_2)$. So (b) holds.

As $g_a(0)<g_a(\kappa)$ in (b), the two cases exclude each other. In (b),
symmetry gives the monotonicity on $(-\infty,0]$. If $g_a$ were unimodal with
mode $m$, then $m\ge0$ would give $g_a(-\kappa)\le g_a(0)$ and $m<0$ would give
$g_a(0)\ge g_a(\kappa)$, while $g_a(\pm \kappa)>g_a(0)$. Finally a local maximum $y$ other
than the points named would make $g_a$ constant on a one-sided
neighborhood of that point. Lemma~\ref{lem:analytic-density} excludes this.
The monotonicities are strict for the same reason, since equal values
at distinct points would give a constant interval between them.

\emph{Step 6, curvature.} The Fourier bound in
Lemma~\ref{lem:analytic-density} gives $g_a(y),g_a'(y)\to0$ as
$y\to\infty$. Evenness and the $\Pi_5$ bound give at most three sign
blocks for $g_a''$ on $(0,\infty)$. Its last sign must be positive. Otherwise $g_a'$ would eventually
be nonincreasing with limit zero, making $g_a$ eventually nondecreasing.
This contradicts $g_a>0$ and $g_a\to0$.
Analyticity and strict monotonicity imply that the first nonconstant
term in the even Taylor series of $g_a$ at zero is negative in case (a)
and positive in case (b). The same signs therefore hold for $g_a''$
near zero. In case (b), $g_a''$ must also be negative somewhere, since
$g_a'$ is positive near zero and negative somewhere beyond $\kappa$.
The positive-half-line sign patterns are consequently $(-,+)$ in
case (a) and $(+,-,+)$ in case (b). Analyticity excludes an interval
of zero curvature and makes the resulting convexities strict. In
case (b), $g_a'(d_1)>0>g_a'(d_2)$, so $d_1<\kappa<d_2$.
\end{proof}

\begin{theorem}[two modes for memory near one]\label{thm:two-modes-near-one}
If $\frac45\le a<1$, that is $0<\varepsilon\le\frac1{10}$, then case (b)
of Theorem~\ref{thm:two-modes} holds. So $g_a$ has exactly two modes.
\end{theorem}

\begin{proof}
In case (a), Khinchine's representation~\cite{dj1988} gives
$\mathcal L^\circ_a\overset{d}=UZ$, with $U$ uniform on $(0,1)$ and
independent of $Z$. Hence Cauchy--Schwarz gives
$\E(\mathcal L^\circ_a)^4\ge\frac95[\E(\mathcal L^\circ_a)^2]^2$.
The moment recursion of~\cite[Theorem~1.4]{glr2026} instead gives
\[
 \frac{\E\mathcal L_a^4}{(\E\mathcal L_a^2)^2}
 =\frac{6(2a^2+2a-1)}{4a-1}\,\operatorname{Beta}(2a,2a),
\]
where $\operatorname{Beta}$ is the beta function. For $a\ge\frac34$, its integral and
$t(1-t)\le\frac14$ give
$\operatorname{Beta}(2a,2a)\le4^{-2a+3/2}\operatorname{Beta}(\frac32,\frac32)=\pi2^{-4a}$.
Put $R(a)=(2a^2+2a-1)/(4a-1)$. On $[\frac34,1]$ we have
$(\log R)'\le2<4\log2$. The first inequality reduces to
$4a(4a^2+a-2)\ge0$. Thus $6\pi R(a)2^{-4a}$ decreases there, and for
$a\ge\frac45$ the displayed moment ratio is at most
\[
 \frac{282\pi}{55\,2^{16/5}}
 <\frac{141}{80}<\frac95,
\]
using $\pi<22/7$ and $2^{1/5}>8/7$. Even moments are unchanged by
symmetrization, so this contradicts case (a). Theorem~\ref{thm:two-modes}
then gives case (b).
\end{proof}

\begin{corollary}[a larger computation-free range]\label{cor:analytic-range}
The conclusion of Theorem~\ref{thm:two-modes-near-one} holds for
$79/100\le a<1$.
\end{corollary}

\begin{proof}
The preceding upper bound $6\pi R(a)2^{-4a}$ decreases for $a\ge3/4$.
At $a=79/100$, use $2^{4/25}>111/100$, which follows from the exact
integer inequality $111^{25}<16\cdot100^{25}$. Substitute this bound
and $\pi<22/7$ to get
\[
 6\pi R(79/100)2^{-79/25}
 <\frac{33517}{18648}<\frac95.
\]
So the kurtosis is again too small for a unimodal density, proving
the claim.
\end{proof}

We locate the transition through the curvature at the origin.
The shape theorem and certified signs leave only a short interval to
check. There the central curvature strictly increases with memory,
so it has at most one zero.

\begin{theorem}[the threshold, computer-assisted]\label{thm:threshold}
Put
\[
 \begin{aligned}
 z_-&=1838418790980431/2748779069440000=0.668812860014597754\ldots,\\
 z_+&=45960469774511/68719476736000=0.668812860014601028\ldots,
 \end{aligned}
\]
so that $z_+-z_-<3.3\cdot10^{-15}$.
\begin{enumerate}[(a)]
\item There is exactly one $a_*\in(\frac12,0.6712]$ with $g_{a_*}''(0)=0$,
and $z_-<a_*<z_+$.
\item For $\frac12<a\le a_*$, case (a) of Theorem~\ref{thm:two-modes}
holds. For $a_*<a<1$, case (b) holds. Thus $a_*$ is the unique transition
from one mode to two, and the critical density is unimodal.
\item We have $g_a''(0)<0$ for $\frac12<a<a_*$ and $g_a''(0)>0$ for
$a_*<a\le0.6712$.
\end{enumerate}
\end{theorem}

Thus the transition is unique even inside the bracket. For
$[79/100,1)$ we also have the moment proof of
Corollary~\ref{cor:analytic-range}, which needs no computation.
For a biased start we lose symmetry, and instead use first differences.

\begin{remark}[the certificate and finite-length thresholds]\label{rem:modes-threshold}
Appendix~\ref{app:threshold} bounds the Fourier integral for $g_a''(0)$
by Taylor models in $a$ and a tail estimate from the fixed-point equation.
It uses scaled moments near $a=\frac12$ and a Hankel test for
$a\ge0.6705$. The certificate takes about 2.8 hours of CPU time.
On the exact finite law,
$D_N(N/2)$ changes sign at $a\approx0.6563$ for $N=4000$ and at
$a\approx0.6592$ for $N=8000$. These two numbers are computed and are not
used in any proof. Proposition~\ref{prop:finite-threshold-location}
proves that, for all sufficiently large even $N$, the zero is unique
on $(0,1)$ and approaches $a_*$ from below, with leading shift
proportional to $N^{-(2a_*-1)}$. It also puts every finite-length
zero below $a_*$, without a length threshold.
\end{remark}


\subsection{Arbitrary initial bias}

Write $h_a$ for the density of $\mathcal L_a$ and
$g_{a,q}(x)=q h_a(x)+(1-q)h_a(-x)$ for initial bias
$\Prob(X_1=1)=q$. Thus $g_{a,1/2}=g_a$.

\begin{theorem}[the interval of bimodal biases]\label{thm:biased-limit-modes}
For every $a\in(1/2,1)$ and $q\in[0,1]$, the density $g_{a,q}$ has at
most two local maxima. There is $b_a\in[0,1/2]$ such that it has exactly
two local maxima if and only if
\[
                 b_a<q<1-b_a.
\]
Outside this open interval, including its endpoints, the density is
unimodal with a unique local maximum. The interval is empty exactly when
$a\le a_*$, in which case $b_a=1/2$.
At finite time, for every $0<\varepsilon<1$, $q\in[0,1]$ and
$N\ge1$, the count law has at most two modal plateaus, including its
endpoints.
\end{theorem}

\begin{proof}
\emph{The mode bound.} Let $P_N$ here denote the count law with first-step bias $q$.
The finite assertion follows from Proposition~\ref{prop:all-bias-modes}.
In particular, the count law cannot have two successive strict valleys.
Neither can $g_{a,q}$. If it did, weak convergence of $S_N/N^a$ and
averages over five short intervals would give five finite-law atoms
with those same forbidden inequalities
(Lemma~\ref{lem:weak-shape-transfer}).
There are also no constant intervals, for any bias, by
Lemma~\ref{lem:analytic-density}. Hence three local maxima $x<y<z$
would give
$x<u<y<v<z$ with
\[
 g_{a,q}(u)<\min\{g_{a,q}(x),g_{a,q}(y)\},\qquad
 g_{a,q}(v)<\min\{g_{a,q}(y),g_{a,q}(z)\}.
\]
For each pair, choose a point in a one-sided maximizing neighborhood
of the smaller maximum. There is a point with strictly smaller value,
since the density is not constant on an interval. The resulting two
valleys contradict the limit bound.

\emph{The bias interval.} Fix $a$, write $f_q=g_{a,q}$ and $h=h_a$,
and call $x<y<z$ a strict valley if $f_q(y)<\min\{f_q(x),f_q(z)\}$.
Proposition~\ref{prop:dominant-half} and the same weak-transfer lemma show that,
for $q\ge1/2$, $f_q$ has no strict valley on $[0,\infty)$ and
$f_q(-x)\le f_q(x)$ for $x\ge0$.
Consequently it has a strict valley if and only if
$f_q(x)>f_q(y)$ for some $x<y\le0$. Such a descent gives the valley
$(x,y,-x)$ by reflection dominance. Conversely, a valley must start at $x<0$. If its middle point $y$
is positive, the no-valley property at $(0,y,z)$ gives
$f_q(0)\le f_q(y)<f_q(x)$.

At $q=1$, conditional unimodality~\cite{glrs2025} excludes a valley,
so $h$ is nondecreasing on $(-\infty,0]$. For fixed $x<y\le0$, the
difference $f_q(x)-f_q(y)$ is affine in $q$ and nonpositive at $q=1$.
If positive at $q\ge1/2$, it remains positive for every $r\in[1/2,q]$.
The valley biases in $[1/2,1]$ are therefore downward closed. They are
relatively open, since the strict inequalities at three fixed points
persist under a small change of bias, and they exclude $1$. Reflection
gives an open interval $(b_a,1-b_a)$, possibly empty.

A strict valley gives two local maxima. Integrability and positivity
supply an outer point on each side below the corresponding witness
height, and we maximize on the two resulting compact intervals. Conversely, two local
maxima give a strict valley because there is no flat interval.
The mode bound excludes a third maximum. A continuous integrable
density without a strict valley is unimodal, and absence of flat
intervals makes its maximum unique. Finally the bias interval is
nonempty exactly when it contains $1/2$, so
Theorem~\ref{thm:threshold} identifies the empty case.
\end{proof}

This proves the initial-bias conjecture
of~\cite[Remark~2.1(b)]{glrs2025}, with its half-width $q(a)=1/2-b_a$.
Trinas and Valle's stronger local limit~\cite[Theorem~3.1]{tv2026},
in an unrefereed arXiv preprint, would also transfer the finite
inequalities. Weak convergence suffices for the proof here.

\subsection{The bias boundary}

To find the endpoint of the bias interval, we compare the slopes
of the conditional density at reflected points. This reduces the
problem to a one-dimensional infimum.

\begin{proposition}[a variational equation for the bias boundary]
\label{prop:bias-boundary}
For $a\in(1/2,1)$ put $U_a=\{t>0:h_a'(t)>0\}$. With the infimum
over an empty set interpreted as $+\infty$,
\begin{equation}\label{eq:bias-boundary}
 b_a=\min\left\{\frac12,\quad
       \inf_{t\in U_a}\frac{h_a'(-t)}{h_a'(t)+h_a'(-t)}\right\}.
\end{equation}
In particular $b_a\to0$ as $a\uparrow1$. If the infimum is attained
at $t>0$ with $h_a'(t)>0$ and $b_a<1/2$, then, with $q=1-b_a$,
\[
 g_{a,q}'(-t)=g_{a,q}''(-t)=0.
\]
At this point the smaller peak disappears.
\end{proposition}

\begin{proof}
The proof of Theorem~\ref{thm:biased-limit-modes} shows that
$h_a'(-t)\ge0$ for $t\ge0$. For $q\ge1/2$ bimodality is equivalent
to a decrease somewhere on the negative half-line, hence to
$g_{a,q}'(-t)<0$ for some $t>0$. Here
\[
 g_{a,q}'(-t)=q h_a'(-t)-(1-q)h_a'(t).
\]
This expression can be negative only if $t\in U_a$, in which case
it is negative exactly when
$1-q>h_a'(-t)/(h_a'(t)+h_a'(-t))$. Taking the union over $t$ gives
\eqref{eq:bias-boundary}, with equality excluded as in the interval
theorem. At an interior minimizer the ratio has derivative zero.
Differentiating its numerator and denominator gives the two displayed
equations. Finally Corollary~\ref{cor:biased-near-one-peaks} puts every
fixed $q\in(0,1)$ inside the interval for all sufficiently large $a$.
Thus $0\le b_a<q$ eventually for every $q>0$, which proves the limit.
\end{proof}

\begin{theorem}[regularity of the bias boundary]\label{thm:bias-continuity}
The function $a\mapsto b_a$ is continuous on $(1/2,1)$, is continuously
differentiable on $(a_*,1)$, and extends
continuously to $b_1=0$. For every $a>a_*$, we have $0<b_a<1/2$,
and the infimum in
\eqref{eq:bias-boundary} has a unique minimizer $t_a\in(0,\kappa_a)$.
This minimizer depends continuously on $a$ throughout $(a_*,1)$.
\end{theorem}

\begin{proof}
For $a>a_*$, the curvature partition in Theorem~\ref{thm:two-modes}
gives $g_a'(t)>0$ on $(0,\kappa_a)$ and $g_a'(t)<0$ on
$(\kappa_a,\infty)$. For $t\ge\kappa_a$ in $U_a$, the ratio in
\eqref{eq:bias-boundary} is therefore at least $1/2$.
Since $b_a<1/2$, any minimizer lies in $(0,\kappa_a)$.
Put
\[
 D_a(t)=h_a'(t)+h_a'(-t),\qquad
 R_a(t)=\frac{g_a'(t)}{D_a(t)}\quad(0<t<\kappa_a).
\]
Lemmas~\ref{lem:balanced-start} and~\ref{lem:positive-conditional-mode}
give $h_a'(-t)>0$ for every $t\ge0$.
We therefore have $D_a>0$, $0<R_a<1/2$, and
$1/2-b_a=\sup_{0<t<\kappa_a}R_a(t)$. Also $R_a(0+)=0$.

The ratio cannot have a strict valley. More precisely,
\begin{equation}\label{eq:bias-quasiconcave}
 R_a(v)\ge\min\{R_a(u),R_a(w)\}
 \qquad(0<u<v<w<\kappa_a).
\end{equation}
Otherwise, choose $s\in(0,1/2)$ strictly between the middle value
and both outer values. With $q=1/2+s$,
$g_{a,q}'(-t)=D_a(t)(s-R_a(t))$ then has signs $-,+,-$
at $-w,-v,-u$, and is positive sufficiently close to zero on the left.
Positivity of the density
and its zero limits at infinity supply a positive derivative farther
left and a negative derivative farther right. There would be three
local maxima, contrary to Theorem~\ref{thm:biased-limit-modes}.

At $\kappa_a$, we have $D_a(\kappa_a)=2h_a'(-\kappa_a)>0$.
Thus $R_a$ vanishes at both ends and attains its positive maximum,
strictly below $1/2$, in the interior. In particular $b_a>0$.
Equation~\eqref{eq:bias-quasiconcave} and analyticity make this
maximum unique, since two maximizers would force a constant interval.
They also make $R_a$ strictly increasing before its maximizer and
strictly decreasing after it.

Lemma~\ref{lem:analytic-density} gives joint continuity of all the
derivatives used here. The signs of $g_a'$ on either side of its
unique positive zero give continuity of $\kappa_a$. Fix $a_0>a_*$
and choose
$0<l<t_{a_0}<r<\kappa_{a_0}$. The strict inequalities
$R_a(l),R_a(r)<R_a(t_{a_0})$ persist near $a_0$.
Equation~\eqref{eq:bias-quasiconcave} bounds every value outside
$[l,r]$ below $R_a(t_{a_0})$, so a maximum on $[l,r]$ is global.
Uniform convergence on this fixed compact
interval proves continuity of the maximum and, by uniqueness, of
its location. This proves continuity above $a_*$. Below it $b_a=1/2$,
Proposition~\ref{prop:bias-cusp} gives continuity at $a_*$, and
Proposition~\ref{prop:bias-boundary} gives the limit at one.

For differentiability, the same three-point argument puts all nearby
maximizers in one compact interval. If $M(a)=1/2-b_a$, maximality gives,
for $s>0$,
\[
 R_{a+s}(t_a)-R_a(t_a)
 \le M(a+s)-M(a)
 \le R_{a+s}(t_{a+s})-R_a(t_{a+s}).
\]
Divide by $s$ and use the continuity of $t_a$ and
$\partial_aR_a$ from Lemma~\ref{lem:analytic-density}. The argument
for $s<0$ is the same. Thus
\begin{equation}\label{eq:bias-envelope-derivative}
 -b_a'=
 \frac{\partial_a h_a'(t_a)\,h_a'(-t_a)
       -h_a'(t_a)\,\partial_a h_a'(-t_a)}
      {\{h_a'(t_a)+h_a'(-t_a)\}^2}.
\end{equation}
The right side is continuous. No derivative of $t_a$ is needed.
\end{proof}

We can use this argument to enclose the boundary numerically, since
three comparisons restrict the search to a bounded interval.
Verified inequalities $R_a(l),R_a(r)<R_a(t)$ with
$0<l<t<r<\kappa_a$ trap the
maximum in $[l,r]$. If $L\le R_a(t)$ and $R_a\le U$ throughout
that interval, then $1/2-U\le b_a\le1/2-L$. If these inequalities hold uniformly, the enclosure covers a whole
parameter interval. More concretely, if
$|R_a(s)-R_{a_0}(s)|\le M|a-a_0|$ for $s=l,t,r$, the two
comparisons at $a_0$ hold throughout $|a-a_0|\le w$ whenever their
gaps exceed $2Mw$. To prove strict decrease of $b_a$ on this interval,
it would suffice to have
\[
 R_{a_2}(s)-R_{a_1}(s)\ge c(a_2-a_1)
 \quad(l\le s\le r,\ a_1<a_2),\qquad c>0.
\]
Indeed, evaluate the later ratio at the earlier maximizer.
This is an interval criterion, not a verified sign bound for the full curve.

\begin{corollary}[a rate for the bias boundary near one]\label{cor:bias-boundary-rate}
Put
\[
 \rho_a=\frac{\Gamma(1+a)^2}{(2a-1)\Gamma(2a)}-1.
\]
If $\rho_a<1/12$, then
\[
 0\le b_a\le\frac{4\rho_a}{1-4\rho_a}.
\]
In particular, as $a\uparrow1$,
$b_a\le8(1-a)+O((1-a)^2)$.
The profile theorem gives the stronger asymptotic bound
\begin{equation}\label{eq:bias-quadratic}
 \limsup_{a\uparrow1}\frac{b_a}{(1-a)^2}
 \le\frac{2}{f_L(\lambda_{\mathrm L})}.
\end{equation}
Thus $b_a=O((1-a)^2)$, and no positive linear lower bound can hold near one.
\end{corollary}

\begin{proof}
The conditional limit has mean $\mu_a=1/\Gamma(1+a)>0$ and variance
$\rho_a\mu_a^2$, by the moment formulas of~\cite[Theorem~1.4]{glr2026}.
Consider the three adjacent intervals of length $\mu_a$ centered at
$-\mu_a$, $0$ and $\mu_a$. Chebyshev's inequality puts at least
$1-4\rho_a$ of the conditional law in the right interval and at most
$4\rho_a$ in the middle. For a start to the left, reflect these
intervals. After mixing, the left interval has mass at least
$(1-q)(1-4\rho_a)$ and the right at least $q(1-4\rho_a)$.
The middle still has mass at most $4\rho_a$.

For $4\rho_a/(1-4\rho_a)<q<1-4\rho_a/(1-4\rho_a)$, the middle
has less mass than either outer interval. All three intervals have
the same length. So, by continuity, we can choose a point in each
with the middle value strictly smaller than the other two.
This rules out unimodality. The two-mode bound in
Theorem~\ref{thm:biased-limit-modes} gives exactly two modes, and
therefore the stated bound on $b_a$. Finally, the logarithmic derivative at $a=1$
of $1+\rho_a$ is $2\psi(2)-2-2\psi(2)=-2$.
Thus $\rho_a=2(1-a)+O((1-a)^2)$, which gives the linear bound.
For the sharper bound, choose $0<u<v$ with $h_a(v)>h_a(u)$ and put
$A=h_a(v)-h_a(u)$ and $B=h_a(-u)-h_a(-v)\ge0$.
The mixture increases from $u$ to $v$ exactly when $q>B/(A+B)$.
For $q\le1/2$, such an increase on the less favored half forces
bimodality. Hence $b_a\le B/(A+B)$ (if this ratio is at least one half,
the bound is automatic). Since $h_a$ is nondecreasing on the negative
half-line, Chebyshev's inequality gives
\[
 u h_a(-u)\le\Prob(\mathcal L_a\le0)\le\rho_a,
 \qquad b_a\le\frac{\rho_a}{u\{h_a(v)-h_a(u)\}}.
\]
Write $a=1-2\varepsilon$ and choose $z_1>z_2>\lambda_{\mathrm L}$
with $f_L(z_1)<f_L(z_2)$.
Take $u=A_\varepsilon-2\varepsilon z_1$ and
$v=A_\varepsilon-2\varepsilon z_2$, using
Theorem~\ref{thm:near-one-profile}. These points tend to one, and
\[
 2\varepsilon\{h_a(v)-h_a(u)\}
 \longrightarrow f_L(z_2)-f_L(z_1)>0.
\]
As $\rho_a\sim4\varepsilon$, division by $(1-a)^2=4\varepsilon^2$
gives a limsup at most $2/\{f_L(z_2)-f_L(z_1)\}$.
Let $z_1\to\infty$ and then $z_2\downarrow\lambda_{\mathrm L}$
to obtain \eqref{eq:bias-quadratic}.
\end{proof}

\subsection{The density at the threshold}

We next ask how far apart the modes are just after the transition.
The needed constants are derivatives at the threshold. Write
\[
 C_* = \left.\partial_a g_a''(0)\right|_{a=a_*},\qquad
 B_*=-g_{a_*}^{(4)}(0),\qquad H_*=h_{a_*}'(0).
\]
Table~\ref{tab:critical-constants} gives strict certified enclosures.
The coefficient $Q_*$ is defined in Proposition~\ref{prop:bias-cusp}.

\begin{center}\small
\captionsetup{hypcap=false}
\begin{tabular}{@{}lrr@{}}
\toprule
Constant & Lower bound & Upper bound\\
\midrule
$C_*$ & $0.6889503945$ & $0.6889503947$\\
$B_*$ & $0.079207784468$ & $0.079207784470$\\
$H_*$ & $0.09707494510086$ & $0.09707494510088$\\
$Q_*$ & $9.86698768$ & $9.86698770$\\
$\sqrt{6C_*/B_*}$ & $7.224131934$ & $7.224131936$\\
\bottomrule
\end{tabular}
\captionof{table}{The constants in the critical expansions. The enclosure
is uniform on a parameter interval containing the certified bracket
for $a_*$ (Appendix~\ref{or:transition-unique}).}
\label{tab:critical-constants}
\end{center}

\begin{proposition}[the critical shape]\label{prop:critical-shape}
The density at the threshold satisfies
\[
 g_{a_*}(u)=g_{a_*}(0)-\frac{B_*}{24}u^4+O(u^6)
 \qquad(u\to0).
\]
For $a>a_*$ let $\kappa_a>0$ be the positive mode of $g_a$. As $a\downarrow a_*$,
\begin{align*}
 \kappa_a&=\sqrt{\frac{6C_*}{B_*}}\sqrt{a-a_*}
                 +o\bigl(\sqrt{a-a_*}\bigr),\\
 g_a(\kappa_a)-g_a(0)&=\frac{3C_*^2}{2B_*}(a-a_*)^2
                 +o\bigl((a-a_*)^2\bigr).
\end{align*}
\end{proposition}

\begin{proof}
On a neighborhood of $a_*$, the certificate bounds
$|\phi_a(t)|$ by $Ct^{-8}$ and $|\partial_a\phi_a(t)|$ by $4t^{-7}$
for large $t$ (Appendix~\ref{or:transition-unique}). On a bounded
frequency interval we instead use the moment series, which converges
locally uniformly, as does its parameter derivative. Together these
bounds let us pass limits and derivatives through the Fourier
integral. We get spatial derivatives through order six continuous
in $a$ and one parameter derivative of the central curvature.
On a smaller parameter interval, the Taylor remainders through order
four are uniform. At criticality, evenness, $g_{a_*}''(0)=0$ and
$g_{a_*}^{(4)}(0)=-B_*<0$ give the first assertion.

Put $\delta=a-a_*>0$. Uniform Taylor expansion gives
\[
 \frac{g_a'(u)}u=C_*\delta+o(\delta)
                  -\frac{B_*+o(1)}6u^2+O(u^4).
\]
For $u=k\sqrt\delta$ this has the sign of $C_*-B_*k^2/6$
when $k$ is fixed away from $\sqrt{6C_*/B_*}$. The unique positive
mode lies between any two such choices with opposite signs, by
Theorem~\ref{thm:two-modes}. Squeezing gives the mode expansion.
Substitute it into
\[
 g_a(u)-g_a(0)=\tfrac12\bigl(C_*\delta+o(\delta)\bigr)u^2
                  -\tfrac1{24}\bigl(B_*+o(1)\bigr)u^4+O(u^6)
\]
to obtain the height expansion.
\end{proof}

\begin{proposition}[how the bias interval opens]\label{prop:bias-cusp}
With the constants of Table~\ref{tab:critical-constants}, put
\[
 Q_* = \frac{\sqrt{2}\,C_*^{3/2}}{3H_*\sqrt{B_*}},
 \qquad Q_*\approx9.86698769.
\]
As $a\downarrow a_*$,
\begin{equation}\label{eq:bias-cusp}
 \frac12-b_a=Q_*(a-a_*)^{3/2}+o\bigl((a-a_*)^{3/2}\bigr).
\end{equation}
On some interval immediately above $a_*$, $b_a$ is strictly decreasing.
For $a>a_*$ sufficiently close to $a_*$, the infimum in
\eqref{eq:bias-boundary} is attained. The minimizer $t_a$ of
Theorem~\ref{thm:bias-continuity} satisfies
\[
 t_a=\sqrt{\frac{2C_*}{B_*}}\sqrt{a-a_*}
       +o\bigl(\sqrt{a-a_*}\bigr)
     =\frac{\kappa_a}{\sqrt3}+o\bigl(\sqrt{a-a_*}\bigr).
\]
At $q=1-b_a$, the stationary point is $-t_a$.
\end{proposition}

\begin{proof}
Proposition~\ref{prop:critical-shape} supplies the required regularity.
The refined certificate in Appendix~\ref{or:transition-unique} encloses
$\pi^{-1}\int_0^\infty t\operatorname{Im}\phi_a(t)\,dt=h_a'(0)$
uniformly on an interval containing $a_*$, giving the positive $H_*$
in Table~\ref{tab:critical-constants}.

Write $\delta=a-a_*$. By \eqref{eq:bias-boundary},
\[
 \frac12-b_a=\max\left\{0,\sup_{t\in U_a}
     \frac{g_a'(t)}{h_a'(t)+h_a'(-t)}\right\}.
\]
The fair-shape theorem makes the numerator nonpositive for $t\ge \kappa_a$.
Proposition~\ref{prop:critical-shape} gives
$\kappa_a=O(\sqrt\delta)$, and $H_*>0$ makes both conditional derivatives
positive on this shrinking interval. Uniform Taylor expansion gives,
for bounded $v\ge0$,
\[
 \frac{g_a'(v\sqrt\delta)}{h_a'(v\sqrt\delta)+h_a'(-v\sqrt\delta)}
 =\frac{\delta^{3/2}}{2H_*}
       \left(C_*v-\frac{B_*v^3}{6}\right)+o(\delta^{3/2}).
\]
The cubic is largest at the unique positive point
$v=\sqrt{2C_*/B_*}$. Maximizing on $[0,\kappa_a/\sqrt\delta]$, whose
right endpoint tends to $\sqrt{6C_*/B_*}$, proves \eqref{eq:bias-cusp}.
The maximum is positive and attained inside $(0,\kappa_a)$, since the ratio
vanishes at both endpoints. Uniform convergence and the unique maximum
of the limiting cubic also force every scaled maximizer to tend to
$\sqrt{2C_*/B_*}$. This proves the assertion about $t_a$.
The bounds on $Q_*$ follow by squaring its formula and inserting the
three rational enclosures for $C_*,B_*$ and $H_*$.

For strict monotonicity, choose a small fixed
$r>0$. For $a$ sufficiently close to $a_*$, the denominator of the
ratio is positive on $[0,r]$ and $\kappa_a<r$. The maximum can therefore
be taken over this fixed compact interval.
The differentiated Fourier bounds justify the expansions
$\partial_a g_a'(t)=C_*t+o(t)$ and
$\partial_a(h_a'(t)+h_a'(-t))=O(1)$, uniformly as
$(a,t)\to(a_*,0)$, while
$g_a'(t)=O((|a-a_*|+t^2)t)$.
The parameter derivative of the ratio is consequently positive for
$0<t\le r$ after shrinking the parameter interval and $r$.
At $a>a_*$ its maximum is positive and occurs at $t>0$.
Evaluating at that maximizer at any larger nearby memory makes the
maximum strictly larger. Hence $b_a$ strictly decreases there.
\end{proof}

This also explains why two different powers appear. Near zero,
the derivative is approximately $C_*(a-a_*)t-B_*t^3/6$.
Its positive part has width of order $\sqrt{a-a_*}$ but height of
order $(a-a_*)^{3/2}$. The modes separate on the width scale.
The bias must overcome the height to remove the smaller peak.

\begin{samepage}
\begin{proposition}[the finite-length threshold]\label{prop:finite-threshold-location}
For every even $N\ge2$, the central curvature has a zero in $(0,1)$.
Every such zero lies below $a_*$. For every sufficiently large
even $N$, there is exactly one zero $a_N$ in the whole interval $(0,1)$.
These thresholds are strictly increasing with $N$ and satisfy
\begin{equation}\label{eq:finite-threshold-shift}
 a_N=a_*-\frac{B_*}{2(2a_*-1)C_*}\,N^{-(2a_*-1)}
          +o\bigl(N^{-(2a_*-1)}\bigr).
\end{equation}
The assertion for large $N$ does not give an effective cutoff or
uniqueness at every smaller even length.
\end{proposition}
\end{samepage}

The proof is in Appendix~\ref{app:finite-threshold}. The main step is
finite. Once the center is a maximum, it remains one at every later
length. This excludes extra zeros away from $a_*$, where the derivative
limit alone would not give a uniform comparison.


\subsection{Density approximation near memory one}

\begin{theorem}[the two peaks near memory one]\label{thm:near-one-profile}
Let $h_a$ denote the density of $\mathcal L_a$, and put
\[
 A_\varepsilon=1+2\varepsilon\log(1/\varepsilon),\qquad
 K_\varepsilon(z)=2\varepsilon
       h_{1-2\varepsilon}(A_\varepsilon-2\varepsilon z),
\]
\[
 H_\varepsilon(z)=4\varepsilon
       g_{1-2\varepsilon}(A_\varepsilon-2\varepsilon z)
       \mathbf1\{z<A_\varepsilon/(2\varepsilon)\}.
\]
Put $\mathcal E_\varepsilon=\varepsilon^{1/4}
\{1+\log(1/\varepsilon)\}^{3/4}$. As $\varepsilon\downarrow0$,
\[
 \|K_\varepsilon-f_L\|_\infty+
 \|H_\varepsilon-f_L\|_\infty=O(\mathcal E_\varepsilon^{1/2}),
 \qquad \|K_\varepsilon-f_L\|_1=O(\mathcal E_\varepsilon^{1/4}).
\]
For every fixed integer $j\ge1$ and $0<\eta<1/2$,
\[
 \|K_\varepsilon^{(j)}-f_L^{(j)}\|_\infty
 =O_{j,\eta}(\mathcal E_\varepsilon^{1/2-\eta}).
\]
If $\kappa_a>0$ is the positive mode of $g_a$, then
\[
 \begin{aligned}
 \kappa_{1-2\varepsilon}
   &=A_\varepsilon-2\varepsilon\lambda_{\mathrm L}
       +O(\varepsilon\mathcal E_\varepsilon^{1/4}),\\
 4\varepsilon g_{1-2\varepsilon}(\kappa_{1-2\varepsilon})
   &=f_L(\lambda_{\mathrm L})+O(\mathcal E_\varepsilon^{1/2}).
 \end{aligned}
\]
For $q\in[0,1]$, define the two-component density
\[
 \widetilde g_{\varepsilon,q}(x)
 =\frac{q}{2\varepsilon}f_L\!\left(\frac{A_\varepsilon-x}{2\varepsilon}\right)
 +\frac{1-q}{2\varepsilon}f_L\!\left(\frac{A_\varepsilon+x}{2\varepsilon}\right).
\]
Uniformly over the initial bias, we have
\[
 \begin{aligned}
 \sup_{q\in[0,1]}2\varepsilon
     \|g_{1-2\varepsilon,q}-\widetilde g_{\varepsilon,q}\|_\infty
       &=O(\mathcal E_\varepsilon^{1/2}),\\
 \sup_{q\in[0,1]}
     \|g_{1-2\varepsilon,q}-\widetilde g_{\varepsilon,q}\|_1
       &=O(\mathcal E_\varepsilon^{1/4}).
 \end{aligned}
\]
\end{theorem}

\begin{proof}
Appendix~\ref{app:fixed-refinements} gives the estimates. The coupling
behind Corollary~\ref{cor:fixed-memory-bridge}, now with a quantitative
comparison of the Poisson jumps, bounds the distribution-function error
by $O(\mathcal E_\varepsilon)$. The fixed-point equation gives uniform
Fourier decay after rescaling, hence bounded spatial derivatives.
Averaging over intervals of length $\mathcal E_\varepsilon^{1/2}$
gives the uniform density bound, and finite differences give the
derivative bounds.
The Landau tail gives the $L^1$ bound, and its quadratic drop from the
maximum gives the mode bound. Reflection and mixing preserve these
estimates, so their constants do not depend on the initial bias.
\end{proof}

\begin{corollary}[both peaks under a biased start]\label{cor:biased-near-one-peaks}
Fix $0<\eta\le1/2$. For all sufficiently small $\varepsilon$, uniformly
over $q\in[\eta,1-\eta]$, the density $g_{1-2\varepsilon,q}$ has exactly
two local maxima $x^-_{\varepsilon,q}<x^+_{\varepsilon,q}$, with
\[
 \begin{aligned}
 x^+_{\varepsilon,q}&=A_\varepsilon-2\varepsilon\lambda_{\mathrm L}
                         +o(\varepsilon),&
 2\varepsilon g_{1-2\varepsilon,q}(x^+_{\varepsilon,q})
                         &=q f_L(\lambda_{\mathrm L})+o(1),\\
 x^-_{\varepsilon,q}&=-A_\varepsilon+2\varepsilon\lambda_{\mathrm L}
                         +o(\varepsilon),&
 2\varepsilon g_{1-2\varepsilon,q}(x^-_{\varepsilon,q})
                         &=(1-q) f_L(\lambda_{\mathrm L})+o(1).
 \end{aligned}
\]
The error terms are uniform for the stated biases.
\end{corollary}

\begin{proof}
Take a small closed interval about $\lambda_{\mathrm L}$ in the
scaled coordinates of the positive peak. The approximation in
Theorem~\ref{thm:near-one-profile} is $qf_L$ there, up to an error
tending uniformly to zero. The reflected component is part of this
error, since its argument $2B_\varepsilon-z$ goes to infinity.
Because $q\ge\eta$, the value at $\lambda_{\mathrm L}$ exceeds
both endpoint values by a gap bounded below uniformly in the bias
once $\varepsilon$ is small enough. Hence the maximum lies inside.
Using $1-q\ge\eta$ gives the corresponding maximum on the negative
side. The intervals are eventually disjoint, and
Theorem~\ref{thm:biased-limit-modes} rules out any further maximum.
Shrinking the fixed intervals proves the position estimates.
The uniform density approximation gives the heights.
\end{proof}

The weights must stay bounded below for this argument. If a bias
tends to zero or one, the approximation remains valid, but its absolute
error may hide the smaller component. For a bound on $b_a$ instead,
see Corollary~\ref{cor:bias-boundary-rate}.
